\section{Comparable prime supply}
\label{app:prime-supply}

The reciprocal provider needs only one fixed-ratio prime band with two-sided
order of magnitude $X/\log X$.  We record an elementary proof.

Define
\[
\theta(x)=\sum_{p\le x}\log p,
\qquad
\psi(x)=\sum_{p^m\le x}\log p
=\log\operatorname{lcm}(1,\ldots,\lfloor x\rfloor).
\]

\begin{lemma}[Central binomial coefficient and the lcm]
For every $n\ge1$,
\[
\binom{2n}{n}
\mid
\operatorname{lcm}(1,\ldots,2n+1).
\]
\end{lemma}

\begin{proof}
Fix a prime $p$.  Legendre's elementary double count gives
\[
v_p(N!)=\sum_{j\ge1}\left\lfloor\frac{N}{p^j}\right\rfloor.
\]
Hence
\[
v_p\binom{2n}{n}
=
\sum_{j\ge1}
\left(
\left\lfloor\frac{2n}{p^j}\right\rfloor
-2\left\lfloor\frac n{p^j}\right\rfloor
\right).
\]
Each summand is $0$ or $1$.  If the sum is $e$, there are $e$ distinct
positive exponents with nonzero summand; the largest is at least $e$ and
satisfies $p^j\le2n$.  Thus $p^e\le2n<2n+1$, so the $p$-adic valuation of the
binomial coefficient is at most that of the displayed lcm.  This holds for
every $p$.
\end{proof}

The central binomial coefficient is maximal in its row, so
\[
4^n
=\sum_{k=0}^{2n}\binom{2n}{k}
\le(2n+1)\binom{2n}{n}.
\]
The preceding lemma therefore yields
\[
\psi(2n+1)
\ge2n\log2-\log(2n+1),
\]
and hence
\[
\psi(x)\ge c_\psi x
\tag{P1}
\]
eventually for some $c_\psi>0$.

Every prime $n<p\le2n$ divides $\binom{2n}{n}$, so
\[
\theta(2n)-\theta(n)
\le\log\binom{2n}{n}
\le2n\log2.
\]
Dyadic telescoping gives
\[
\theta(x)\le C_\theta x
\tag{P2}
\]
for an absolute $C_\theta$; for example $C_\theta=4\log2$ works.

Since
\[
\psi(x)=\sum_{m\ge1}\theta(x^{1/m}),
\]
(P2) gives
\[
0\le\psi(x)-\theta(x)
\le
C_\theta
\sum_{2\le m\le\log_2x}x^{1/m}
=O(\sqrt x\log x)
=o(x).
\]
Combining with (P1), there are constants $c_\theta,C_\theta>0$ such that
\[
c_\theta x\le\theta(x)\le C_\theta x
\tag{P3}
\]
eventually.

Choose an integer $\Lambda>1$ with
\[
C_\theta<c_\theta\Lambda.
\]
Then for large $X$,
\[
\theta(\Lambda X)-\theta(X)
\ge(c_\theta\Lambda-C_\theta)X.
\]
Since every prime in $(X,\Lambda X]$ contributes at most
$\log(\Lambda X)$, this gives the lower bound
\[
\#\{p:X<p\le\Lambda X\}
\gg\frac X{\log X}.
\]
For the upper bound, every such prime contributes at least $\log X$, while
\[
\theta(\Lambda X)\le C_\theta\Lambda X.
\]
Thus:

\begin{proposition}[Comparable prime band]\label{prop:prime-band}
For some fixed $\Lambda>1$ and constants $c,C>0$,
\[
c\frac X{\log X}
\le
\#\{p:X<p\le\Lambda X\}
\le
C\frac X{\log X}
\]
for every sufficiently large $X$.
\end{proposition}

No numerical choice of $\Lambda,c,C$ survives the provider boundary.


\section{The prime-order restricted-fold theorem}
\label{app:restricted-fold}

We prove the additive theorem used in Theorem~\ref{thm:restricted-fold}.
Let $S$ be an additive group of prime order $p$, let
\[
A\hookrightarrow S,
\qquad |A|=r,
\]
and let
\[
\Sigma_h:\operatorname{FinSub}_h(A)\to S
\]
be the $h$-element sum map.  We prove the stronger bound
\[
|\im\Sigma_h|
\ge
\min\{p,h(r-h)+1\}.
\tag{D1}
\]

Choose a coordinate $S\cong\mathbf F_p$ only inside this proof.

\subsection{A finite-grid coefficient detector}

\begin{lemma}\label{lem:grid-detector}
Let $F$ be a field. Suppose
\[
P\in F[x_1,\ldots,x_h]
\]
has total degree at most $d_1+\cdots+d_h$ and the coefficient of
\[
x_1^{d_1}\cdots x_h^{d_h}
\]
is nonzero. If finite sets $A_i\subset F$ satisfy $|A_i|>d_i$, then $P$ is
nonzero at some point of $A_1\times\cdots\times A_h$.
\end{lemma}

\begin{proof}
Choose $B_i\subseteq A_i$ of cardinality $d_i+1$. For $a\in B_i$ put
\[
\ell_{i,a}(t)
=
\prod_{\substack{b\in B_i\\b\ne a}}
\frac{t-b}{a-b}.
\]
Lagrange interpolation implies that the linear functional
\[
L_i(f)
=
\sum_{a\in B_i}
\frac{f(a)}{\prod_{b\in B_i,\ b\ne a}(a-b)}
\]
satisfies
\[
L_i(t^j)=0\quad(j<d_i),
\qquad
L_i(t^{d_i})=1.
\]
Because the total degree of $P$ is at most $\sum_i d_i$, the tensor
$L_1\otimes\cdots\otimes L_h$ kills every monomial of $P$ except the target
monomial and returns precisely its coefficient. If $P$ vanished on the whole
grid, the same tensor, being a linear combination of grid evaluations, would
return zero, a contradiction.
\end{proof}

\subsection{The Vandermonde coefficient}

Let
\[
\operatorname{Vdm}(x)
=
\prod_{i<j}(x_j-x_i).
\]
For distinct nonnegative integers $k_1,\ldots,k_h$, set
\[
m=\sum_i k_i-\frac{h(h-1)}2.
\]
Then
\[
[x_1^{k_1}\cdots x_h^{k_h}]
(x_1+\cdots+x_h)^m\operatorname{Vdm}(x)
=
m!\,
\frac{\prod_{i<j}(k_j-k_i)}{\prod_i k_i!}.
\tag{D2}
\]
Indeed, expand
\[
\operatorname{Vdm}(x)
=
\det(x_i^{j-1})
=
\sum_{\sigma\in S_h}
\operatorname{sgn}(\sigma)
\prod_i x_i^{\sigma(i)-1}.
\]
The multinomial theorem transforms the desired coefficient into
\[
\frac{m!}{\prod_i k_i!}
\det\bigl((k_i)_{j-1}\bigr),
\]
where $(k)_j$ is the falling factorial. Since $(t)_{j-1}$ is monic of degree
$j-1$, its evaluation determinant equals the ordinary Vandermonde determinant,
which is $\prod_{i<j}(k_j-k_i)$. This proves (D2).

\subsection{All admissible exponent sums}

For every
\[
0\le m\le h(r-h)
\]
there are distinct integers
\[
0\le k_1<\cdots<k_h\le r-1
\]
with
\[
\sum_i k_i
=m+\frac{h(h-1)}2.
\tag{D3}
\]
To see this, put $R=r-h$ and write $m=qh+s$ with $0\le s<h$. If $q<R$, take
$h-s$ of the numbers $\lambda_i$ equal to $q$ and the remaining $s$ equal to
$q+1$; if $q=R$, necessarily $s=0$ and take them all equal to $R$. Then
\[
0\le\lambda_1\le\cdots\le\lambda_h\le R,
\qquad
\sum_i\lambda_i=m,
\]
and
\[
k_i=(i-1)+\lambda_i
\]
have the required properties.

\subsection{Image growth}

Put
\[
M=\min\{p-1,h(r-h)\}.
\]
Suppose for contradiction that
\[
|\im\Sigma_h|\le M.
\]
Enlarge the image to a set $C\subset\mathbf F_p$ of cardinality $M$ and set
\[
P(x_1,\ldots,x_h)
=
\prod_{c\in C}
\left(\sum_i x_i-c\right)
\operatorname{Vdm}(x).
\]
This polynomial vanishes on $A^h$: if two coordinates coincide, the
Vandermonde vanishes; if they are distinct, their sum lies in $\im\Sigma_h\subseteq C$.

Use (D3) with $m=M$ to choose distinct
\[
0\le k_1<\cdots<k_h\le r-1
\]
whose sum has the required degree. The top homogeneous part of $P$ is
\[
(\sum_i x_i)^M\operatorname{Vdm}(x),
\]
so (D2) says that the coefficient of
$x_1^{k_1}\cdots x_h^{k_h}$ equals
\[
M!\,
\frac{\prod_{i<j}(k_j-k_i)}{\prod_i k_i!}.
\]
Here $M\le p-1$ and $k_i\le r-1\le p-1$, so all factorials are nonzero in
$\mathbf F_p$; the $k_i$ are distinct, so every difference is nonzero as well.
Thus the coefficient is nonzero. Lemma~\ref{lem:grid-detector} contradicts the
grid vanishing. Therefore (D1) holds.

Finally, the statement is coordinate independent. An isomorphism between two
prime-order additive groups induces a bijection of their $h$-element finite
subobjects and commutes with the sum map:
\[
\begin{tikzcd}
\operatorname{FinSub}_h(A) \arrow[r,"\Sigma_h"] \arrow[d,"\operatorname{FinSub}_h(u)"']
& S \arrow[d,"u"]\\
\operatorname{FinSub}_h(uA) \arrow[r,"\Sigma_h"']
& S'.
\end{tikzcd}
\]
Hence image cardinality is intrinsic. In particular, if
\[
h(r-h)+1\ge p,
\]
then (D1) gives $|\im\Sigma_h|=p$, so $\Sigma_h$ is surjective. This proves
Theorem~\ref{thm:restricted-fold}.
